\documentclass[aspectratio=169, 11pt]{beamer}

% ========== Theme ==========
\usetheme{Madrid}
\usecolortheme{whale}
\setbeamertemplate{navigation symbols}{}
\setbeamertemplate{footline}[frame number]

% ========== Packages ==========
\usepackage{tikz}
\usetikzlibrary{arrows.meta, positioning, shapes.geometric, calc, fit, backgrounds}
\usepackage{booktabs}
\usepackage{graphicx}
\usepackage{amsmath}
\usepackage{xcolor}

% ========== Colours ==========
\definecolor{highlight}{RGB}{220, 70, 70}
\definecolor{completed}{RGB}{80, 185, 100}
\definecolor{pending}{RGB}{180, 180, 180}
\definecolor{inprogress}{RGB}{90, 140, 240}

% ========== Title ==========
\title{Weekly Progress Report}
\subtitle{Week 31 --- Draft Feedback Review and Revision Plan}
\author{Student ID: 11314389}
\institute{BCI Control System Project}
\date{April 22, 2026}

\begin{document}

% ===================================================
% TITLE PAGE
% ===================================================
\begin{frame}
    \titlepage
\end{frame}

% ===================================================
% SLIDE 1: Agenda
% ===================================================
\begin{frame}[shrink=7]{Week 31 Agenda --- Feedback Discussion on the Final Draft}

    \begin{block}{Purpose of This Meeting}
        Use the annotated comments on the report draft to identify the main problems in the current version and confirm how the next revision should be prioritised.
    \end{block}

    \vspace{2mm}
    \begin{exampleblock}{Topics for Discussion}
        \begin{enumerate}
            \item Main issues raised in the annotated feedback on the draft
            \item Why the current \textbf{Methodology} chapter is not yet working well enough
            \item Specific changes planned for the next revision
            \item Which remaining sections should be prioritised after the \textbf{Methods} rewrite
        \end{enumerate}
    \end{exampleblock}

    \vspace{2mm}
    \begin{alertblock}{Main Goal}
        Leave this meeting with a clear and defensible revision plan for the next draft.
    \end{alertblock}
\end{frame}

% ===================================================
% SLIDE 2: Core Feedback
% ===================================================
\begin{frame}[shrink=7]{1. Core Issues Identified in the Feedback}

    \begin{columns}[T]
        \column{0.49\textwidth}
        \textbf{Main Writing Problems}
        \begin{itemize}
            \item Methods and results are mixed together
            \item Some explanation appears too interpretive for a methods chapter
            \item Several claims are stronger than the evidence currently shown
            \item Some statements need references or weaker wording
        \end{itemize}

        \vspace{2mm}
        \textbf{Main Structural Problems}
        \begin{itemize}
            \item The boundary between Literature Review and Methods is not always clear
            \item Some technical content appears before it is properly introduced
            \item The chapter does not yet read like a clean reproducible protocol
        \end{itemize}

        \column{0.47\textwidth}
        \begin{alertblock}{Most Important Overall Message}
            The current chapter is not being judged mainly on whether the technical method is good or bad. The bigger problem is that it does not yet read like a proper \textbf{Methods} chapter.
        \end{alertblock}

        \vspace{2mm}
        \begin{exampleblock}{Interpretation}
            The next revision should focus first on clarity, reproducibility, and chapter boundaries, rather than on adding more technical novelty.
        \end{exampleblock}
    \end{columns}
\end{frame}

% ===================================================
% SLIDE 3: Methodology Problems
% ===================================================
\begin{frame}[shrink=7]{2. Why the Current Methodology Chapter Needs Revision}

    \begin{columns}[T]
        \column{0.5\textwidth}
        \textbf{What Is Missing or Unclear}
        \begin{itemize}
            \item Dataset description is incomplete
            \item Sampling rates and channel context need to be clearer
            \item Preprocessing choices need stronger justification
            \item Normalisation procedure is not explicit enough
            \item Model details are not yet reproducible enough
            \item Training / validation / testing protocol needs to be clearer
        \end{itemize}

        \column{0.47\textwidth}
        \textbf{What Is Currently in the Wrong Place}
        \begin{itemize}
            \item Accuracy improvements appearing inside the methods chapter
            \item Comparative performance statements before the results chapter
            \item Interpretation of why performance changed before the evidence is shown
        \end{itemize}

        \vspace{2mm}
        \begin{alertblock}{Expected Revision Direction}
            Chapter 3 should answer ``what was done and how'', while Chapter 4 should answer ``what happened'' and ``what it means''.
        \end{alertblock}
    \end{columns}
\end{frame}

% ===================================================
% SLIDE 4: Planned Methods Rewrite
% ===================================================
\begin{frame}[shrink=7]{3. Planned Rewrite of the Methods Chapter}

    \begin{table}
        \centering
        \small
        \begin{tabular}{p{4.2cm}p{6.3cm}c}
            \toprule
            \textbf{Feedback Area} & \textbf{Planned Change} & \textbf{Priority} \\
            \midrule
            Results mixed into Methods
            & Move all accuracy gains, comparisons, and outcome statements into Chapter 4
            & \textcolor{highlight}{High} \\
            Dataset description too thin
            & Add sampling rate, channel count, class definition, and protocol table
            & \textcolor{highlight}{High} \\
            Preprocessing not reproducible enough
            & State filter range, epoch length, normalisation rule, and ICA status explicitly
            & \textcolor{highlight}{High} \\
            Model description too abstract
            & Add clearer EEGTransformer architecture explanation, figure, and config summary
            & \textcolor{highlight}{High} \\
            Training protocol unclear
            & Clarify subject-dependent, LOSO, pooled training, and fine-tuning procedures
            & \textcolor{highlight}{High} \\
            Some claims too strong
            & Add references where needed or reduce certainty in wording
            & \textcolor{inprogress}{Medium} \\
            \bottomrule
        \end{tabular}
    \end{table}
\end{frame}

% ===================================================
% SLIDE 5: Proposed New Chapter Logic
% ===================================================
\begin{frame}[shrink=6]{4. Proposed Logic for the Next Draft}

    \begin{columns}[T]
        \column{0.48\textwidth}
        \textbf{Chapter 2: Literature Review}
        \begin{itemize}
            \item Explain the existing MI classification landscape
            \item Introduce benchmark datasets and prior methods
            \item Set up why the project approach is reasonable
        \end{itemize}

        \vspace{2mm}
        \textbf{Chapter 3: Methods}
        \begin{itemize}
            \item Datasets used
            \item Preprocessing pipeline
            \item EEGTransformer architecture
            \item Training protocol
            \item Channel reduction protocol
            \item RL formulation
        \end{itemize}

        \column{0.48\textwidth}
        \textbf{Chapter 4: Results and Discussion}
        \begin{itemize}
            \item Classification results
            \item Channel reduction results
            \item RL performance
            \item ICA / filter ablation results
            \item Interpretation and limitations
        \end{itemize}

        \vspace{2mm}
        \begin{exampleblock}{Benefit of This Separation}
            The report becomes easier to follow because the marker no longer has to mentally separate protocol, evidence, and interpretation while reading the same section.
        \end{exampleblock}
    \end{columns}
\end{frame}

% ===================================================
% SLIDE 6: Wider Draft Improvements
% ===================================================
\begin{frame}[shrink=7]{5. Additional Improvements Planned Beyond Methods}

    \begin{columns}[T]
        \column{0.5\textwidth}
        \textbf{Writing / Presentation}
        \begin{itemize}
            \item Reduce over-strong interpretive language
            \item Use more cautious phrasing where evidence is limited
            \item Improve figure readability, especially overview diagrams
            \item Standardise terms such as \textit{participant} where appropriate
        \end{itemize}

        \column{0.46\textwidth}
        \textbf{Final Report Structure}
        \begin{itemize}
            \item Check alignment with handbook and report template
            \item Ensure Project Management is presented correctly
            \item Improve appendix organisation
            \item Include repository evidence and supporting materials clearly
        \end{itemize}

        \vspace{2mm}
        \begin{alertblock}{Reason for These Changes}
            Even if the technical work is strong, weak structure and unclear academic writing could reduce the overall impression of the report.
        \end{alertblock}
    \end{columns}
\end{frame}

% ===================================================
% SLIDE 7: Revision Plan
% ===================================================
\begin{frame}[shrink=6]{6. Revision Plan After This Meeting}

    \begin{table}
        \centering
        \small
        \begin{tabular}{p{5.0cm}p{5.4cm}c}
            \toprule
            \textbf{Task} & \textbf{Planned Outcome} & \textbf{Order} \\
            \midrule
            Rewrite Methods chapter
            & Clean reproducible protocol chapter with no embedded results
            & 1 \\
            Reorganise Results / Discussion
            & Move outcomes and interpretation to the correct sections
            & 2 \\
            Tighten language and citations
            & Fewer unsupported claims and clearer academic tone
            & 3 \\
            Check format against template
            & Draft looks closer to a final submission
            & 4 \\
            Final consistency pass
            & Stable terminology, figure quality, and appendix structure
            & 5 \\
            \bottomrule
        \end{tabular}
    \end{table}

    \vspace{2mm}
    \begin{alertblock}{Question for This Meeting}
        After the Methods rewrite, which section should be treated as the next highest priority: \textbf{Discussion}, \textbf{overall structure / formatting}, or \textbf{appendix / submission packaging}?
    \end{alertblock}
\end{frame}

\end{document}
